\documentclass[sigconf]{acmart}
\usepackage{tikz}
\usetikzlibrary{positioning, fit, calc, backgrounds, shapes.geometric}
\usepackage{xcolor}

% Revision highlight command — renders in blue for easy identification
% Remove this command and all \rev{} wrappers before final submission
\newcommand{\rev}[1]{{\color{blue}#1}}

% \usepackage{graphicx}
% \usepackage{booktabs}
% \usepackage{amsmath}
% \usepackage{microtype} 

\AtBeginDocument{%
  \providecommand\BibTeX{{%
    \normalfont B\kern-0.5em{\scshape i\kern-0.25em b}\kern-0.8em\TeX}}}

\setcopyright{acmcopyright}
\copyrightyear{2026}
\acmYear{2026}
\acmDOI{XXXXXXX.XXXXXXX}
\acmConference[COLIEE '26]{Competition on Legal Information Extraction and Entailment}{June 12, 2026}{Singapore}

\begin{document}

\title{Strategic Formulation of the Decoupled Architecture for Legal Case Entailment: Overcoming the Context Paradox and Lexical Traps}

\author{Somdev Ganguli}
\affiliation{\institution{Jadavpur University} \city{Kolkata} \country{India}}
\email{somdevganguli@gmail.com}

\author{Vibhan Dutta}
\affiliation{\institution{Jadavpur University} \city{Kolkata} \country{India}}
\email{bivandatta07@gmail.com}

\author{Amit Barman}
\affiliation{\institution{Jadavpur University} \city{Kolkata} \country{India}}
\email{amitbarman811@gmail.com}

\author{Debasis Ganguly}
\affiliation{\institution{University of Glasgow} \city{Glasgow} \country{UK}}
\email{debasis.ganguly@glasgow.ac.uk}

\author{Sudip Kumar Naskar}
\affiliation{\institution{Jadavpur University} \city{Kolkata} \country{India}}
\email{sudip.naskar@gmail.com}

\renewcommand{\shortauthors}{JUNLLP Team}

\begin{abstract}
Legal case entailment involves the challenge of navigating dense, standardized judicial texts to identify whether a specific historical paragraph logically supports a new legal decision fragment. Recent methodologies for COLIEE Task 2 have frequently utilized large-scale generative models; however, these approaches can be computationally intensive and remain susceptible to domain-specific lexical biases. This paper explores the hypothesis that effective legal entailment benefits from precise semantic boundary detection as a specialized alternative to generative approximation. \rev{To motivate and justify our architectural decisions, we introduce two diagnostic metrics: Truncation Damage Index (TDI), which quantifies structural information loss due to context limits, and Lexical Trap Divergence (LTD), which identifies instances of ``lexical collapse'' where model confidence is driven by surface-level overlap rather than genuine entailment. These diagnostics reveal the specific failure modes that our architecture is designed to address.} To address these identified bottlenecks, we evaluate a three-stage decoupled architecture consisting of a hybrid retrieval phase using Weighted Reciprocal Rank Fusion (wRRF), a LegalBERT cross-encoder utilizing Pooling-based Multi-Head Attention (PMA), and a three-gate max-gap algorithm to manage predictive variance. Our integrated pipeline achieves a leaderboard F1-score of 0.3942\rev{, establishing a high-precision operating point (Precision: 0.6721) that prioritizes reliable entailment identification over aggressive recall}. Furthermore, experimental results suggest that the LTD \rev{framework} may serve as a useful \rev{diagnostic signal for identifying susceptibility to lexical collapse in} high-stakes legal reasoning tasks.
\end{abstract}

\begin{CCSXML}
<ccs2012>
   <concept>
       <concept_id>10010147.10010178.10010179</concept_id>
       <concept_desc>Computing methodologies~Natural language processing</concept_desc>
       <concept_significance>500</concept_significance>
   </concept>
   <concept>
       <concept_id>10002951.10003317</concept_id>
       <concept_desc>Information systems~Information retrieval</concept_desc>
       <concept_significance>500</concept_significance>
   </concept>
</ccs2012>
\end{CCSXML}

\ccsdesc[500]{Computing methodologies~Natural language processing}
\ccsdesc[500]{Information systems~Information retrieval}

\keywords{Legal Case Entailment, Context Paradox, Lexical Traps, Decoupled Architecture, Natural Language Processing}

\maketitle

\begin{figure*}[t]
  \centering
  \includegraphics[width=\textwidth, keepaspectratio]{archite_diag.pdf}
  \caption{System Description of the Decoupled Architecture for COLIEE 2026 Task 2. Stage 1 utilizes a Weighted Reciprocal Rank Fusion (wRRF) to maximize initial recall. Stage 2 employs a Context-Safe LegalBERT Cross-Encoder, isolating the query to prevent truncation while utilizing a 4-head PMA layer and focal loss to learn deep semantic boundaries. Stage 3 applies a 3-Gate Safe Max-Gap algorithm to eliminate predictive variance.}
  \label{fig:architecture}
\end{figure*}

\section{Introduction}

The increasing availability of digital legal documents has driven the need for specialized Natural Language Processing (NLP) models capable of navigating the complexities of the legal domain \cite{ariai2025nlp}. Within this field, Legal Case Entailment (CL-DE) remains a particularly challenging task. As formally defined by Task 2 of the Competition on Legal Information Extraction and Entailment (COLIEE) \cite{rabelo2024overview}, the objective is to predict the decision of a new case by identifying entailment from previous relevant cases. Specifically, given a query containing a decision fragment from a new case and a corresponding noticed case, the system must automatically identify the exact paragraph within the noticed case that entails the query. Crucially, in this context, the ``decision'' does not refer to the final case verdict, but rather to a specific conclusion expressed by the judge (the entailed fragment).

Recent high-performing architectures for this task have increasingly relied on parameter scaling, advanced dense embeddings, and complex generative pipelines \cite{li2025kalm, li2026lexentail}. For example, the \rev{NOWJ team's benchmark of 0.3195 F1 \cite{nguyen2024nowj} on the COLIEE 2025 Task 2 dataset} utilized extensive ensembling of large language models. Our preliminary experiments with Small Language Models (SLMs) yielded similar insights: while prompt-engineered SLMs achieve high recall by capturing broad semantic relationships, they often suffer from lower precision. Because generative models are primarily optimized for probabilistic token prediction rather than strict boundary detection, they can be easily misled by domain-specific legal phrasing \cite{mishra2025investigating}. In contrast, bidirectional encoders like the BERT family function as discriminative classifiers \cite{nogueira2019passage}, offering the structural capacity to establish the firm semantic boundaries required for high-precision logical filtering. However, standard encoder methods still face limitations when processing lengthy legal texts.

In this paper, we argue that effective, fully automated legal entailment requires prioritizing precise semantic extraction over generative approximation. We focus on addressing two specific challenges that frequently limit performance in this domain:
\begin{itemize}
    \item \textbf{The Context Paradox (see Section 2.3):} Standard transformer architectures operate with strict token limits. We empirically observe that prepending excessive base-case context to queries often forces the model to truncate the actual ground-truth evidence required to prove entailment.
    \item \textbf{The Lexical Trap (see Section 2.2):} Judicial opinions frequently utilize standardized phrasing and boilerplate statutes. Traditional retrieval metrics over-index on this surface-level vocabulary overlap, assigning high relevance scores to paragraphs that lack true logical entailment.
\end{itemize}

To resolve these dual challenges, we introduce a novel, 3-stage Decoupled Architecture. This pipeline leverages a hybrid retrieval mechanism to cast a wide, diverse recall net \cite{kiyohara2025off}, followed by a specialized cross-encoder designed to evaluate entailing logic in a strictly context-safe environment. Finally, a deterministic gating algorithm eliminates late-stage ranking variance. By focusing on isolating true semantic logic and actively defending against lexical overlap biases, our framework achieves an official leaderboard F1-score of 0.3942, offering a robust standard for legal natural language inference.

\rev{Our system was submitted under the team name \textbf{JUNLLP}. The official submission utilized Run ID \textbf{submission (1)}, corresponding to the 3-seed ensemble configuration described in Section 4.2.}

\section{Diagnostic Framework}

Before formulating our pipeline, we mathematically quantify the specific dataset characteristics that cause standard models to fail. Existing generation and entailment metrics often struggle to capture the nuances of open-ended legal reasoning because they retain inherent biases toward surface-level lexical overlap \cite{mccoy2019right}. To provide a rigorous foundation for our architectural decisions, we define two custom diagnostic metrics: Lexical Trap Divergence (LTD) and the Truncation Damage Index (TDI).

\begin{table*}[t]
  \caption{Empirical Demonstration of the Lexical Trap. Standard retrieval ($S_{lex}$) fails by assigning high scores to Hard Negatives due to boilerplate overlap. Our fine-tuned Cross-Encoder ($S_{sem}$) successfully suppresses these traps while highly scoring the True Positives,effectively reducing the Lexical Trap Divergence (LTD).}
  \label{tab:ltd_evidence}
  \small
  \begin{tabular}{p{1.5cm} p{10cm} c c c}
    \toprule
    \textbf{Query ID} & \textbf{Paragraph Type \& Snippet} & \textbf{$S_{lex}$} & \textbf{$S_{sem}$} & \textbf{LTD}\\
    \midrule
    928 & \textit{Hard Neg:} [14] The applicant's submissions with respect to the visa officer's credibility are without merit be... & 1.000 & 0.089 & 2.416 \\
         & \textit{True Pos:} [12] This Court has made it clear that the determination of an applicant's occupational qualificatio... & 0.835 & 0.431 & 0.552 \\
    \midrule
    930 & \textit{Hard Neg:} [9] The Board dismissed Mr. Akhter's appeal of the decision refusing his wife and daughter's applica... & 0.507 & 0.125 & 0.709 \\
         & \textit{True Pos:} [20] The two issues in this application are concerned with the interpretation of the paragraph 117(9... & 0.119 & 0.185 & -0.052 \\
    \midrule
    931 & \textit{Hard Neg:} [23] The applicant argued that the CRA officer should have waived all penalties and interests, not o... & 1.000 & 0.130 & 2.037 \\
         & \textit{True Pos:} [18] As already mentioned, the Act and its regulations are silent as to what criteria are to be used... & 0.404 & 0.235 & 0.220 \\
    \bottomrule
  \end{tabular}
\end{table*}

\subsection{Structuring the Lexical Trap Divergence}

The ``Lexical Trap'' is an adversarial phenomenon where models substitute deep semantic reasoning for shallow, keyword-based matching. This is uniquely pervasive in legal texts due to the heavy reliance on boilerplate language and standardized statutory phrasing. 

To quantify the exact threshold at which a model falls into this trap, we formulate the Lexical Trap Divergence (LTD) metric. The LTD serves as a domain-specific quantification of the ``Lexical Overlap Bias'' previously identified in general Natural Language Inference tasks \cite{mccoy2019right}.

To derive the LTD, we establish two distinct probability distributions for a given query fragment ($q$) and a candidate paragraph ($p$). First, we define the sparse Lexical Similarity ($S_{lex}$) using a MinMax-normalized BM25 scoring function to bound the vocabulary overlap between $0$ and $1$:
\begin{equation}
    S_{lex}(q, p) = \frac{\text{BM25}(q, p) - \text{BM25}_{min}}{\text{BM25}_{max} - \text{BM25}_{min}}
\end{equation}

Second, we define the Semantic Entailment confidence ($S_{sem}$). Establishing strict logical entailment requires the continuous logit output of our deeply fine-tuned LegalBERT cross-encoder. We isolate the sigmoidal probability of the entailment class:
\begin{equation}
    S_{sem}(q, p) = \sigma(\mathbf{W} \cdot \mathbf{h}_{[\text{CLS}]} + b)
\end{equation}

Drawing inspiration from the Kullback-Leibler (KL) divergence \cite{kullback1951information}, the LTD calculates the information penalty incurred when high lexical overlap does not correlate with actual semantic entailment:
\begin{equation}
    \mathrm{LTD}(q, p) = S_{lex}(q, p) \cdot \log \left( \frac{S_{lex}(q, p)}{S_{sem}(q, p) + \epsilon} \right)
\end{equation}
where $\epsilon$ is a smoothing constant to prevent undefined logarithms.

\subsection{Empirical Evidence of the Lexical Trap}

To demonstrate the prevalence of this issue within the COLIEE dataset, we isolated adversarial examples from our validation fold, presented in Table~\ref{tab:ltd_evidence}. These ``Hard Negatives'' recite general jurisprudential standards, sharing near-identical keywords with the query fragment. Consequently, standard sparse retrieval algorithms assign them an artificially high $S_{lex}$ score.

Statistical analysis across our validation set confirms that paragraphs exhibiting an LTD score greater than 1.5 correlate with a 74\% false-positive rate in standard BM25 retrieval, empirically validating LTD as a robust diagnostic signal for identified lexical traps. Conversely, the ground-truth entailed paragraphs---the ``True Positives''---often utilize specific, localized language that shares fewer exact tokens with the query, resulting in lower sparse retrieval scores. In our training split, an unweighted reliance on sparse retrieval resulted in a 52\% miss rate of True Positives at Rank-1. However, as demonstrated in Table~\ref{tab:ltd_evidence}, our optimized cross-encoder successfully recognizes the lack of logical affirmation in the Hard Negatives (suppressing $S_{sem}$),effectively reducing the LTD penalty and suggesting that discriminative boundary detection is critical for legal entailment.

\rev{We emphasize that LTD is employed as a post-hoc diagnostic signal rather than an active component of our training or inference pipeline. Its primary role is to empirically validate the architectural choices described in Section 3---specifically, the need for a deeply fine-tuned cross-encoder capable of suppressing lexical traps that surface-level retrieval cannot resolve.}

% [A2] LTD role clarification added per Reviewer 1 feedback

\subsection{The Context Paradox and Truncation Damage Index}

The second primary bottleneck is structural. Standard bidirectional transformer architectures impose a hard limit on input length. In long-document retrieval tasks, researchers frequently attempt to enrich the neural network by prepending the broader surrounding context of the base case to the query \cite{lewis2020retrieval}.

We identify this approach as the ``Context Paradox.'' Bounded token environments dictate that attempting to enrich the model's understanding by feeding it broader context directly reduces its capacity to evaluate the target data. 

\rev{To quantify the severity of this paradox, we introduce the Truncation Damage Index (TDI). Conceptually, let $f(x)$ denote the probability density function over token position $x$, representing the likelihood that a ground-truth entailing paragraph begins at position $x$ within the tokenized base case. In the continuous formulation, the TDI is defined as the cumulative probability mass falling strictly beyond the transformer's maximum context window $L_{\max}$:}
\begin{equation}
    \rev{TDI = \int_{L_{\max}}^{\infty} f(x)\,dx}
\end{equation}

\rev{Since legal documents are discrete token sequences, we estimate $f(x)$ empirically as follows. We iterate over all 925 training queries in the COLIEE 2026 dataset. For each query, we tokenize the full base case document using the exact tokenizer employed in our pipeline (\texttt{nlpaueb/legal-bert-base-uncased}). We then locate each ground-truth entailing paragraph within the tokenized base case via substring matching at the token level, recording the starting token index $x_i$ of each matched paragraph. This yields a collection of $N$ starting positions $\{x_1, x_2, \ldots, x_N\}$ across all queries. The empirical density function is therefore:}
\begin{equation}
    \rev{\hat{f}(x) = \frac{1}{N} \sum_{i=1}^{N} \delta(x - x_i)}
\end{equation}
\rev{where $\delta(\cdot)$ is the Dirac delta function. Substituting this discrete approximation, the TDI reduces to a simple counting ratio:}
\begin{equation}
    \rev{TDI = \frac{|\{x_i : x_i > L_{\max}\}|}{N} \approx 0.1007}
\end{equation}

\rev{In plain terms: we count how many ground-truth entailing paragraphs begin beyond the 512-token boundary and divide by the total number of ground-truth paragraphs examined.}

Our empirical analysis of the COLIEE data revealed a TDI of \rev{approximately} 10.07\% (analyzed further in Section \rev{5.1}). This indicates that for over one in ten positive cases, the neural network is \rev{structurally prevented} from predicting the correct answer because the vital tokens are arbitrarily deleted prior to inference. 

Furthermore, our formulation of the TDI codifies the structural information loss historically observed as the ``Lost in the Middle'' attention degradation \cite{liu2024lost}. Even when tokens are theoretically retained, monolithic long-context models suffer from severe context degradation. This diagnostic metric provides the empirical justification for the ``Decoupling'' strategy implemented in Stage 2 of our pipeline.

% [A3] Added f(x) derivation with full empirical methodology and fixed section reference

\section{The Decoupled Architecture}

To actively neutralize the Lexical Trap and the Context Paradox formalized in Section 2, we propose a 3-stage Decoupled Architecture. The fundamental philosophy of this pipeline is strict operational isolation: we utilize contextually enriched queries to maximize initial recall, but aggressively strip that context prior to deep semantic classification to preserve structural integrity.

\subsection{Stage 1: Enriched Query wRRF Retrieval}

The objective of Stage 1 is to cast a wide, high-recall net over the noticed case to extract the top-$k$ candidate paragraphs ($k=80$, empirically determined to maximize recall while maintaining cross-encoder throughput). To maximize the retrieval surface area, we construct an Enriched Query ($Q_{enriched}$) by concatenating the raw entailed fragment ($F_{query}$) with a compressed semantic context ($C_{base}$) extracted from the broader base case using localized lexical scoring:
\begin{equation}
    Q_{enriched} = C_{base} \oplus F_{query}
\end{equation}

\rev{Specifically, we segment the base case text into sentences, score each sentence against the query fragment $F_{query}$ using BM25 (via the \texttt{rank\_bm25} library), and select the top-$k=2$ highest-scoring sentences to form the compressed context $C_{base}$. This targeted extraction ensures the enriched query captures the most topically relevant fragments of the base case without introducing excessive noise.}

% [A7] Context localized BM25 clarification added per Reviewer 2 feedback

This enriched query is passed through two parallel retrieval branches. The sparse lexical branch utilizes BM25 to capture exact statutory phrasing, while the dense semantic branch utilizes the \texttt{BAAI/\allowbreak bge-m3} embedding model \cite{chen2024bge} to capture broad contextual alignments. To merge these disparate vector spaces, we apply a Weighted Reciprocal Rank Fusion (wRRF) \cite{cormack2009reciprocal}:
\begin{equation}
    wRRF(d) = \sum_{m \in \{bge, bm25\}} \frac{w_m}{\kappa + r_m(d)}
\end{equation}
where $r_m(d)$ is the rank of document $d$ in model $m$, and $\kappa=60$ is a standard mathematical constant established in the original RRF literature \cite{cormack2009reciprocal}. This heavily asymmetric weighting strategy ($w_{bge}=1.0$, $w_{bm25}=0.2$) is functionally essential in the legal domain to prevent lexical metrics from overwhelming true semantic intent. By relegating the sparse retriever to a supplementary exact-match safety net, our wRRF calibration ensures dense embeddings remain the primary retrieval driver.

\subsection{Stage 2: Context-Safe PMA Cross-Encoder}

Stage 2 represents the critical ``Decoupling'' phase of our architecture. To explicitly circumvent the Truncation Damage Index (TDI) established in Section 2.3, we completely discard the base case context $C_{base}$ used in Stage 1. Only the raw query fragment $F_{query}$ and the retrieved candidate paragraphs $P_{candidate}$ are passed to the cross-encoder. 

Our core semantic classifier is built upon \texttt{nlpaueb/\allowbreak legal-bert-\allowbreak base-uncased} \cite{chalkidis2020legal}. However, standard BERT architectures rely exclusively on the output of the \texttt{[CLS]} token for classification \cite{reimers2019sentence}. We empirically observed that this single-vector bottleneck restricts the model's capacity to aggregate logical conditions scattered across long legal paragraphs. To resolve this, we adapt the Pooling by Multihead Attention (PMA) mechanism, originally introduced for permutation-invariant set processing \cite{lee2019set}, to dynamically route features across the document sequence. Unlike static Mean or Max pooling, the learnable queries in the PMA layer allow the architecture to dynamically suppress recurrent legal boilerplate while amplifying sparse entailing signals across lengthy sequences.

Given the sequence of hidden states $\mathbf{H} \in \mathbb{R}^{L \times d}$ extracted from the final transformer layer, we project the keys and values via learnable weight matrices $\mathbf{W}_K, \mathbf{W}_V \in \mathbb{R}^{d \times d}$:
\begin{equation}
    \mathbf{K} = \mathbf{H}\mathbf{W}_K, \quad \mathbf{V} = \mathbf{H}\mathbf{W}_V
\end{equation}

Instead of deriving queries from the text sequence itself, we introduce a static, learnable parameter matrix $\mathbf{Q}_{pool} \in \mathbb{R}^{h \times d}$ serving as the independent pooling queries for $h=4$ attention heads (striking an optimal balance between capturing sufficient semantic diversity and avoiding latency bloat). The raw attention scores $\mathbf{S} \in \mathbb{R}^{h \times L}$ are computed as:
\begin{equation}
    \mathbf{S} = \frac{\mathbf{Q}_{pool} \mathbf{K}^T}{\sqrt{d}}
\end{equation}

After applying the sequence padding mask, we compute the softmax across the sequence dimension to derive the attention weights. These weights are multiplied by the values to extract the head-specific pooled representations $\mathbf{P} \in \mathbb{R}^{h \times d}$:
\begin{equation}
    \mathbf{P} = \text{softmax}(\mathbf{S})\mathbf{V}
\end{equation}

Finally, the matrix $\mathbf{P}$ is flattened into a single semantic vector $\mathbf{v}_{out} \in \mathbb{R}^{h \cdot d}$, which is passed into the final linear classifier. This mechanism enables the architecture to anchor its attention on multiple distant semantic boundaries within the text without being constrained by the standard \texttt{[CLS]} bottleneck.

\subsubsection{Focal Loss Optimization for Hard Negatives}
Training an encoder to safely navigate Lexical Traps requires carefully curated negative sampling. We utilize the wRRF ranks from Stage 1 to bucket our negative samples, exposing the model to a strict ratio of 6 Hard Negatives (Lexical Traps) : 4 Semi-Hard : 2 Easy per True Positive. This specific curriculum forces the model to construct strict decision boundaries around Lexical Traps without forgetting basic entailment features. 

To prevent the overwhelming presence of difficult samples from destabilizing the gradient during training, we optimize the network using Focal Loss \cite{lin2017focal}:
\begin{equation}
    \mathrm{FL}(p_t) = -\alpha_t (1 - p_t)^\gamma \log(p_t)
\end{equation}
where $p_t$ is the predicted probability for the ground-truth class. We utilize hyperparameters $\alpha=0.75$ and $\gamma=2.0$, enforcing a focal-style reweighting mechanism that heavily penalizes the model when it confidently predicts a lexical trap as an entailment.

\subsection{Stage 3: 3-Gate Safe Max-Gap Variance Reduction}

Standard cross-encoders frequently suffer from predictive variance, outputting clustered probability distributions where the True Positive is mathematically indistinguishable from adjacent noise. To stabilize the final decision boundary, we implement a deterministic 3-Gate Safe Max-Gap algorithm. 

As illustrated in our system architecture (Figure \ref{fig:architecture}), this mechanism evaluates the sorted list of candidate probabilities $P = [p_1, p_2, \dots, p_n]$ and their corresponding paragraphs $C = [c_1, c_2, \dots, c_n]$. To determine the optimal dynamic threshold, the algorithm scans the top 8 candidates (as our empirical analysis indicates ground-truth entailment rarely exists beyond the 8th candidate) to identify the maximum probability gap $\Delta_{max} = p_k - p_{k+1}$, establishing $k$ as the candidate cut-off index. We define the decay ratio at this boundary as $\rho_k = p_k / (p_{k+1} + \epsilon)$. 

To consolidate our rigorous acceptance criteria, we isolate the boolean acceptance condition $\mathcal{C}_{acc}$ as:
\begin{equation}
    \mathcal{C}_{acc} \equiv (p_k \geq 0.35) \land (\Delta_{max} \geq 0.06) \land (\rho_k \geq 1.30)
\end{equation}

These gating constants were determined through a localized grid search on the validation fold and remained frozen during test-set evaluation to ensure algorithmic generalizability.

The final prediction function $\mathcal{G}(P, C)$ is then formalized as a piecewise cascade rejection operation:
\begin{equation}
    \mathcal{G}(P, C) =
    \begin{cases}
        \{c_1\} & \text{if } p_1 < 0.10 \\
        \{c_1\} & \text{if } p_1 < 0.20 \land (p_1 - p_2) < 0.03 \\
        \bigcup_{i=1}^{k} \{c_i\} & \text{if } \mathcal{C}_{acc} \\
        \{c_1\} & \text{otherwise}
    \end{cases}
\end{equation}

The logic governing these gates operates as a defensive cascade:
\begin{itemize}
    \item \textbf{Gate 1 (Noise Filter):} The first condition recognizes a failure of entailment logic, forcing the model to gracefully halt and return only the top candidate as a low-confidence fallback.
    \item \textbf{Gate 2 (Plateau Guard):} The second condition identifies low-confidence plateaus where the model is effectively guessing blindly among the top candidates.
    \item \textbf{Gate 3 (Rigorous Acceptance):} The $\mathcal{C}_{acc}$ condition executes a multi-paragraph acceptance only when a mathematically significant, high-confidence boundary is detected.
\end{itemize}
This mechanism effectively neutralizes late-stage variance, prioritizing high-precision safety over reckless multi-paragraph generation.

\section{Experiments and Results}

\subsection{Dataset and Evaluation Metrics}

We evaluate our Decoupled Architecture on the official dataset provided by the Competition on Legal Information Extraction and Entailment (COLIEE) 2026 for Task 2\footnote{\url{https://coliee.org/COLIEE2026/tasks/task2}}. The objective is to match a specific decision fragment from the base case to its logically entailing paragraph in the noticed case. 

Unlike standard reading comprehension benchmarks, COLIEE Task 2 does not feature an independent human annotator baseline, as the ground-truth entailment pairs are derived directly from the historical judicial record. The detailed distribution of queries and candidate paragraphs across the training and test splits is provided in Table \ref{tab:dataset_dist}.

\begin{table}[h]
  \caption{Distribution of candidate paragraphs and true positive entailment depths across the training and test splits.}
  \label{tab:dataset_dist}
  \begin{tabular}{@{}l c c@{}}
    \toprule
    \textbf{Metric} & \textbf{Training Split} & \textbf{Test Split} \\
    \midrule
    Total Queries & 925 & 100 \\
    Total Candidates & 32,717 & 4,673 \\
    Avg Candidates / Query & 35.37 & 46.73 \\
    Max Candidates & 315 & 289 \\
    Avg True Positives / Query & 1.28 & 2.94 \\
    Single-Label Queries & 78.81\% & 61.00\% \\
  \bottomrule
\end{tabular}
\end{table}

Following the official COLIEE evaluation framework \cite{rabelo2024overview}, system performance is measured using \rev{micro-averaged} Precision ($P$), Recall ($R$), and F1-measure. Given the strict nature of the task, F1 is the primary ranking metric.

% [B2] FIXED: Changed from "macro-averaged" to "micro-averaged" to match the actual code implementation.

\subsection{Implementation Details}

To ensure reproducibility, we detail the configurations of our pipeline. All experiments were conducted utilizing a single NVIDIA Tesla T4 GPU in a constrained cloud environment. 

In Stage 1, text tokenization and sparse retrieval were executed using the \texttt{rank\_bm25} library, while dense embeddings were generated using the 1024-dimensional \texttt{BAAI/\allowbreak bge-m3} model \cite{chen2024bge}. For Stage 2, our custom PMA Cross-Encoder was initialized with weights from \texttt{nlpaueb/\allowbreak legal-bert-\allowbreak base-uncased} \cite{chalkidis2020legal}. The network was optimized using AdamW \cite{loshchilov2017decoupled} for 4 epochs with a per-device batch size of 8 and 2 gradient accumulation steps. We utilized a learning rate of $2\text{e-}5$, a weight decay of $0.01$, and a warmup ratio of $0.1$. The maximum sequence length was strictly bounded to 512 tokens, with context aggressively pruned to solely include the query fragment and candidate paragraph. 

To further stabilize our predictive boundaries, our final submission utilized a 3-seed ensemble (random seeds 42, 43, and 44), averaging the output probabilities prior to the gating mechanism.

\subsection{Baseline Models}

We benchmark our Decoupled Architecture against three distinct classes of retrieval and entailment systems:
\begin{itemize}
    \item \textbf{Lexical Baselines:} Standard BM25 \cite{robertson2009probabilistic}. We evaluate a Rank-1 ($k=1$) BM25 retrieval, utilizing standard term frequency saturation parameters ($k_1 = 1.5, b = 0.75$).
    \item \textbf{Dense Retrieval Baselines:} Zero-shot dense retrieval using \texttt{BAAI/\allowbreak bge-m3} \cite{chen2024bge}, demonstrating the performance of modern embeddings without domain-specific entailment fine-tuning.
    \item \textbf{Historical Reference:} The \rev{NOWJ 2025} system \cite{nguyen2024nowj}, which achieved an F1 of 0.3195 utilizing complex LLM ensembling. \rev{The NOWJ scores are reproduced from \cite{nguyen2024nowj} on the COLIEE 2025 Task 2 dataset. We cite this as a historical reference point from a prior competition year; differences in test set composition preclude direct comparison.}
\end{itemize}

\begin{table}[t]
  \caption{Performance comparison on the COLIEE Task 2 Test Split. Our 3-stage Decoupled Architecture outperforms both lexical/dense baselines and prior SOTA LLM ensembles.}
  \label{tab:main_results}
  \begin{tabular}{@{}l c c c@{}}
    \toprule
    \textbf{Model / Architecture} & \textbf{Precision} & \textbf{Recall} & \textbf{F1-Score} \\
    \midrule
    \textit{Lexical \& Dense Baselines} & & & \\
    BM25 (Rank-1) & 0.5500 & 0.1871 & 0.2792 \\
    BGE-m3 (Zero-Shot) & 0.5000 & 0.1701 & 0.2538 \\
    \midrule
    \textit{Historical Reference (COLIEE 2025)} & & & \\
    \rev{NOWJ (2025)} \cite{nguyen2024nowj} & 0.3788 & 0.2762 & 0.3195 \\
    \midrule
    \textit{Proposed Framework} & & & \\
    Vanilla BERT (Standard \texttt{[CLS]}) & \rev{0.6400} & \rev{0.2177} & \rev{0.3249} \\
    LegalBERT (Standard \texttt{[CLS]}) & \rev{0.6800} & \rev{0.2313} & \rev{0.3452} \\
    \textbf{JUNLLP (Ours: PMA + Ensemble)} & \textbf{0.6721} & \textbf{0.2789} & \textbf{0.3942} \\
  \bottomrule
\end{tabular}
\end{table}

\rev{Fine-tuning uses random weight initialization and our submission averages three seeds (Section~4.2); metrics therefore vary slightly across runs. The scores in Table~\ref{tab:main_results} for the \texttt{[CLS]} baselines match the cumulative ablation ladder (Table~\ref{tab:component_ablation}, Rows~1--2) on the same test split. Small differences between those headline numbers and auxiliary tables---e.g., thresholding experiments (Table~\ref{tab:threshold_ablation}) evaluated on a fixed single-seed checkpoint for controlled comparison---fall within the seed-to-seed variance we observe and are not contradictory.}

\begin{figure*}[t]
  \centering
  % The \hspace*{1.5cm} acts as an invisible block pushing the image to the right.
  % Tweak the 1.5cm value (e.g., 0.5cm or 2.0cm) until it is visually perfect!
  \makebox[\textwidth][l]{\hspace*{1.67cm}\includegraphics[width=0.95\textwidth]{attention_heatmap_massive_font.pdf}}
  \caption{Attention Rollout demonstrating the PMA cross-encoder successfully anchoring on critical legal terminology while ignoring stop-words, illustrating robust boundary detection.}
  \label{fig:attention_heatmap}
\end{figure*}

\subsection{Results and Analysis}

The comparative evaluation results are presented in Table \ref{tab:main_results}. Our proposed Decoupled Architecture (JUNLLP) achieves a high-precision operating point for legal entailment, with an official leaderboard F1-score of 0.3942, outperforming standard retrieval baselines and the historical NOWJ reference. As depicted in the Precision-Recall curve (Figure \ref{fig:pr_curve}), our architecture maintains robust boundary detection, sharply contrasting with the baseline dense retrieval collapse on hard negatives.

\begin{figure}[t]
  \centering
  \includegraphics[width=\linewidth]{test_pr_curve.png}
  \caption{Precision-Recall (PR) Curve contrasting the robust boundary detection of our PMA architecture against the baseline dense retrieval collapse on hard negatives.}
  \label{fig:pr_curve}
\end{figure}

The low performance of the BM25 baseline (F1 = 0.2792) provides strong empirical support for our hypothesis regarding the Lexical Trap; relying solely on keyword overlap in the legal domain results in a severe precision drop due to boilerplate statutes. While broad generative ensembles frequently attempt to inflate overall F1 scores by over-predicting candidates, our system enforces a strict logical boundary. In the legal domain, precision is critical; a system that confidently surfaces incorrect case law as logical entailment poses significant risks to jurisprudential review. 

Furthermore, the ablation between the standard \texttt{[CLS]}-pooled LegalBERT and our PMA-upgraded architecture highlights the necessity of multi-head semantic routing for lengthy legal texts. Visual evidence of this robust boundary detection is provided in the attention rollout (Figure \ref{fig:attention_heatmap}), which demonstrates the PMA cross-encoder successfully anchoring on critical legal terminology. By safely discarding the context block (neutralizing the TDI), leveraging the 4-head PMA, and utilizing the 3-seed ensemble to stabilize the 3-Gate algorithm, our system successfully suppresses false positives without sacrificing initial recall. This is perfectly illustrated in Figure \ref{fig:pareto_frontier}, demonstrating that our architecture pushes the latency-precision Pareto frontier significantly further than generative models or fast baselines alone.

\section{Ablation Studies and Discussion}

To justify the architectural restraint of our proposed pipeline, we conducted targeted ablation studies evaluating alternative paradigms that ultimately proved detrimental in the legal entailment domain.

\rev{\subsection{Component-Wise Ablation}

To rigorously isolate the contribution of each architectural component, we conducted a six-configuration ablation study on the test data (100 queries). Each row in Table~\ref{tab:component_ablation} introduces exactly one additional component over the previous configuration, enabling a clear attribution of performance gains. All configurations share the same Stage 1 retrieval (wRRF top-80) and identical hyperparameters; only the specified component is toggled.

\begin{table}[h]
  \caption{Component-wise ablation on the test set (100 queries). Each row adds one component to the previous configuration. Metrics are micro-averaged.}
  \label{tab:component_ablation}
  \begin{tabular}{@{}l c c c@{}}
    \toprule
    \textbf{Configuration} & \textbf{Prec.} & \textbf{Recall} & \textbf{F1} \\
    \midrule
    1.\;Vanilla BERT [CLS] & 0.6400 & 0.2177 & 0.3249 \\
    2.\;LegalBERT [CLS] & 0.6800 & 0.2313 & 0.3452 \\
    3.\;LegalBERT + PMA & 0.6900 & 0.2347 & 0.3503 \\
    4.\;LegalBERT + PMA + Focal & 0.7100 & 0.2415 & 0.3604 \\
    5.\;+ 3-Seed Ensemble (Top-1) & 0.7200 & 0.2449 & 0.3655 \\
    \textbf{6.\;Full Pipeline (3-Gate)} & \textbf{0.7232} & \textbf{0.2755} & \textbf{0.3990} \\
    \bottomrule
  \end{tabular}
\end{table}

Several observations emerge from Table~\ref{tab:component_ablation}. First, switching from vanilla \texttt{bert-base-uncased} (Row~1, F1 $0.3249$) to the domain-adapted \texttt{legal-bert-base-uncased} (Row~2, F1 $0.3452$) yields a clear absolute gain of $+0.0203$ in F1, confirming that legal-domain pre-training improves entailment reasoning under identical training. Second, stacking the 4-head PMA layer on LegalBERT (Row~3, F1 $0.3503$) already improves over Row~2; adding Focal Loss with bucketed hard-negative curriculum (Row~4, F1 $0.3604$) compounds this, showing that PMA and lexical-trap-aware training are complementary. Third, the 3-seed probability ensemble with Top-1 decisions (Row~5, F1 $0.3655$) stabilizes rankings before gating. Finally, the full 3-Gate Max-Gap pipeline (Row~6) attains the strongest configuration (F1 $0.3990$), driven primarily by a recall increase from $0.2449$ to $0.2755$ alongside a modest precision lift from $0.7200$ to $0.7232$, which empirically justifies dynamic multi-paragraph acceptance over a fixed Top-1 rule.}

\rev{\subsection{Negative Sampling Ratio Sensitivity}

The bucketed negative sampling strategy described in Section~3.2.1 employs a ratio of 6 Hard : 4 Semi-Hard : 2 Easy negatives per True Positive. To validate this specific ratio, we trained the full LegalBERT + PMA + Focal configuration under four different Hard:Semi-Hard:Easy ratios, each maintaining 12 total negatives per positive for a controlled comparison. Results on the validation fold are presented in Table~\ref{tab:neg_ratio}.

\begin{table}[h]
  \caption{Negative sampling ratio ablation on the validation fold. All configurations use LegalBERT + PMA + Focal with Top-1 thresholding and identical hyperparameters.}
  \label{tab:neg_ratio}
  \begin{tabular}{@{}l c c c@{}}
    \toprule
    \textbf{Ratio (H:SH:E)} & \textbf{Prec.} & \textbf{Recall} & \textbf{F1} \\
    \midrule
    \textbf{6:4:2 (Ours)} & \textbf{0.6486} & \textbf{0.2920} & \textbf{0.4027} \\
    4:4:4 (Uniform) & 0.6270 & 0.2822 & 0.3893 \\
    8:2:2 (Hard-heavy) & 0.6324 & 0.2847 & 0.3926 \\
    2:4:6 (Easy-heavy) & 0.6270 & 0.2822 & 0.3893 \\
    \bottomrule
  \end{tabular}
\end{table}

The results confirm that our 6:4:2 ratio achieves the highest F1 (0.4027). Notably, the Uniform (4:4:4) and Easy-heavy (2:4:6) ratios converge to identical performance (F1 = 0.3893, P = 0.6270, R = 0.2822), revealing a performance floor: once hard-negative exposure falls below a critical proportion, the model's discriminative capacity saturates at the same baseline regardless of easy-negative volume. The Hard-heavy (8:2:2) ratio recovers partially (F1 = 0.3926) but remains below our chosen ratio, suggesting that excessive hard-negative exposure without sufficient gradient-stabilizing easy examples introduces training variance. These results indicate that the 6:4:2 ratio occupies a narrow effective range between insufficient and excessive hard-negative pressure.}

\subsection{Evaluating Truncation Risk in Concatenated Retrieval}

A common baseline optimization in document retrieval involves concatenating multiple highly-ranked candidates into a single cross-encoder prompt to reduce inference latency. To evaluate the feasibility of this approach within the COLIEE Task 2 framework, we profiled the exact token distributions of our dataset using the \texttt{legal-bert-base-uncased}\cite{chalkidis2020legal} tokenizer. 

Our analysis revealed an average query fragment length of 44.6 tokens and an average candidate paragraph length of 140.1 tokens, with the 90th percentile of candidates reaching 261.0 tokens. Given LegalBERT's strict 512-token limit (which includes 3 overhead special tokens), the average available budget for candidate evaluation is strictly 464.4 tokens. Even when evaluated independently (decoupled), 2.52\% of exceptionally long paragraphs naturally exceed this boundary, triggering strict Truncation Damage (TDI). 

Attempting to concatenate merely the top three wRRF candidates ($k=3$) produces an average sequence length of 467.9 tokens (3 + 44.6 + [140.1 $\times$ 3]), guaranteeing context overflow in the average case. If any candidate within that concatenated sequence falls into the 90th percentile (261.0 tokens), the sequence catastrophically overflows, actively deleting the required ground-truth evidence prior to inference. Therefore, processing all $k=80$ candidates completely independently is structurally mandatory to preserve semantic integrity, justifying the linear $\mathcal{O}(N)$ computational overhead.

\subsection{Generative vs. Discriminative Prompting}

To isolate the efficacy of our architecture's discriminative boundary detection, we ablated our pipeline against a generative Small Language Model (SLM) baseline utilizing Qwen2.5-7B-Instruct (4-bit quantized). We established a two-stage pipeline: retrieving the top-4 candidates via BGE-m3, and prompting the SLM to independently verify strict legal entailment. 

\rev{We note that the SLM baseline operates on a reduced candidate set (top-4 from BGE-m3) compared to our primary pipeline (top-80 from wRRF), reflecting the computational constraints of autoregressive inference. While this introduces an asymmetry in retrieval quality, our primary claim is that discriminative boundary detection outperforms generative reasoning at the per-candidate level---as evidenced by the precision gap between our full cross-encoder system (0.6721; Table~\ref{tab:main_results}) and the best generative strategy (0.5752, CoT). A controlled comparison at equal retrieval depth is left for future work.}

% [A4] SLM fairness acknowledgment added per Reviewer 1 feedback

\begin{figure}[t]
  \centering
  \includegraphics[width=\linewidth]{pareto_frontier.pdf}
  \caption{Latency vs. Performance Pareto Frontier. The plot visually confirms that the Decoupled Architecture operates at an optimal intersection of computational efficiency and entailment precision.}
  \label{fig:pareto_frontier}
\end{figure}

\begin{table}[h]
  \caption{Ablation of Qwen2.5-7B Prompting Strategies (Top-4 Candidates). CoT achieves the highest precision among generative methods, but all remain strictly inferior to the cross-encoder.}
  \label{tab:slm_ablation}
  \begin{tabular}{@{}l c c c@{}}
    \toprule
    \textbf{Prompting Strategy} & \textbf{Precision} & \textbf{Recall} & \textbf{F1-Score} \\
    \midrule
    Zero-Shot & 0.5273 & 0.1973 & 0.2871 \\
    Few-Shot (Standard) & 0.5038 & \textbf{0.2245} & 0.3106 \\
    Chain-of-Thought (CoT) & \textbf{0.5752} & 0.2211 & \textbf{0.3194} \\
    Legal IRAC Framework & 0.5366 & \textbf{0.2245} & 0.3165 \\
  \bottomrule
\end{tabular}
\end{table}

We evaluated four distinct prompting paradigms, detailed in Table \ref{tab:slm_ablation}. Standard Chain-of-Thought (CoT) prompting achieved the highest overall generative performance (F1: 0.3194, Precision: 0.5752) by forcing explicit logical reasoning. Interestingly, the domain-specific IRAC (Issue, Rule, Application, Conclusion) framework tied for the highest recall (0.2245) but suffered a slight drop in precision. This suggests that while formal legal frameworks aid in capturing potential matches, they occasionally overcomplicate straightforward entailments when interacting with the Lexical Trap, leading to false positives.

Crucially, \rev{none of these generative methods reach the precision of our Context-Safe PMA cross-encoder pipeline (0.6721; Table~\ref{tab:main_results}).} To mathematically quantify this failure, we adapted our Lexical Trap Divergence (LTD) metric for the generative SLM. By extracting the raw token log-probabilities for the model's final entailment classification token (e.g., ``YES''), we isolate the generative semantic confidence, $S_{llm}$. The adapted generative LTD is formulated as:
\begin{equation}
    LTD_{slm}(q, p) = S_{lex}(q, p) \cdot \log \left( \frac{S_{lex}(q, p)}{S_{llm}(q, p) + \epsilon} \right)
\end{equation}

Our empirical analysis of $LTD_{slm}$ revealed a clear difference in behavior. When the SLM successfully resisted a Lexical Trap (True Negatives), it achieved an average LTD of 1.6638, decoupling its semantic reasoning from the baseline keyword overlap. However, when the model hallucinated the entailments in the boilerplate statutes (False Positives), mirroring retrieval failures observed in broader LLM search agents \cite{sun2023chatgpt}, the average LTD plunged to -0.2696. This negative divergence provides empirical evidence of ``Lexical Collapse''---the point at which the generative model abandons deep semantic reasoning, and its confidence artificially spikes to match the shallow lexical overlap of the Hard Negative. By replacing this vulnerable generative reasoning with a strictly bounded discriminative hyperplane, our architecture achieves superior entailment mapping.

\subsection{Context Concatenation vs. Decoupling}
We evaluated a `coupled' variant of our architecture, where the base case context ($C_{base}$) from Stage 1 was retained and prepended to the query during Stage 2 cross-encoding. As predicted by our Truncation Damage Index (TDI) framework, this resulted in a sharp degradation in performance, plunging the F1-score to 0.2335 (Precision: 0.4600, Recall: 0.1565).

The prepended context consumed the majority of the 512-token transformer window. Furthermore, even when tokens were not explicitly truncated, the attention weights suffered from the well-documented ``Lost in the Middle'' phenomenon \cite{liu2024lost}, where standard transformers fail to utilize vital information buried deep within extended contexts. This ablation empirically supports the need for our ‘Decoupling’ mechanism.

\subsection{\rev{Thresholding Dynamics: Precision--Recall Trade-off Analysis}}
Standard classification pipelines default to thresholding based on a static confidence value (e.g., $\sigma > 0.5$) or a naive \textit{argmax}. To demonstrate why this fails in legal entailment, we tracked the predictive variance across the test set candidates.

\begin{figure}[h]
  \centering
  \includegraphics[width=\linewidth]{3gatesafemaxgap.pdf}
  \caption{Analysis of predictive variance. The plot illustrates how standard static thresholding leads to false-positive spikes, whereas our 3-Gate Max-Gap algorithm cuts off retrieval at the mathematical plateau.}
  \label{fig:3gate_plot}
\end{figure}

\begin{table*}[t]
  \caption{\rev{Comprehensive threshold grid search on the validation fold (925 queries). All configurations use the same single-seed LegalBERT + PMA + Focal checkpoint. F1\textsubscript{micro} is the official COLIEE ranking metric; F1\textsubscript{macro} measures per-query consistency.}}
  \label{tab:threshold_ablation}
  \small
  \begin{tabular}{@{}l c c c c c@{}}
    \toprule
    \textbf{Strategy} & \textbf{$\tau$} & \textbf{Precision} & \textbf{Recall} & \textbf{F1\textsubscript{micro}} & \textbf{F1\textsubscript{macro}} \\
    \midrule
    Static (0.10) & 0.10 & 0.1166 & 0.4331 & 0.1837 & 0.2598 \\
    Static (0.15) & 0.15 & 0.1849 & 0.3990 & 0.2527 & 0.3428 \\
    Static (0.20) & 0.20 & 0.2551 & 0.3674 & 0.3011 & 0.4072 \\
    Static (0.25) & 0.25 & 0.3145 & 0.3528 & 0.3326 & 0.4546 \\
    Static (0.30) & 0.30 & 0.3658 & 0.3382 & 0.3515 & 0.4789 \\
    Static (0.35) & 0.35 & 0.4149 & 0.3260 & 0.3651 & 0.5101 \\
    Static (0.40) & 0.40 & 0.4662 & 0.3187 & 0.3786 & 0.5289 \\
    Static (0.45) & 0.45 & 0.4924 & 0.3139 & 0.3834 & 0.5410 \\
    Static (0.50) & 0.50 & 0.5292 & 0.3090 & \textbf{0.3902} & 0.5560 \\
    Static (0.55) & 0.55 & 0.5381 & 0.2920 & 0.3785 & 0.5438 \\
    Static (0.60) & 0.60 & 0.5545 & 0.2847 & 0.3762 & 0.5427 \\
    Static (0.65) & 0.65 & 0.5659 & 0.2822 & 0.3766 & 0.5472 \\
    Static (0.70) & 0.70 & 0.5876 & 0.2774 & 0.3769 & 0.5535 \\
    \midrule
    Top-1 (Argmax) & --- & 0.6000 & 0.2701 & 0.3725 & 0.5544 \\
    \textbf{3-Gate Max-Gap (Ours)} & \textbf{dyn.} & \textbf{0.5228} & \textbf{0.3066} & 0.3865 & \textbf{0.5579} \\
    \bottomrule
  \end{tabular}
\end{table*}

\rev{As shown in Figure \ref{fig:3gate_plot} and detailed in Table \ref{tab:threshold_ablation}, the static threshold at $\tau = 0.50$ achieves the highest micro-F1 (0.3902), marginally above our 3-Gate method (0.3865). However, this aggregate metric does not capture per-query reliability. Our 3-Gate mechanism attains the highest macro-F1 across all 15 configurations (0.5579), indicating superior per-query consistency. Micro-F1 is dominated by high-candidate queries where aggressive recall is rewarded; macro-F1, by weighting each query equally, better reflects reliability across the full case distribution. The 3-Gate method achieves this robustness because it dynamically adapts its acceptance boundary to the localized probability landscape of each individual query, rather than applying a corpus-level constant. For legal applications where consistent per-query reliability outweighs aggregate corpus-level optimization, this operating point represents the more principled choice.}

% [A1] 3-Gate narrative fully reframed per Reviewer 1 feedback
% [B1] Threshold grid search data incorporated from notebook Cell 9

\subsection{Qualitative Human Error Analysis}
Because COLIEE Task 2 ground truth is derived from historical judicial citations, the human standard is theoretically bounded at 1.0. Therefore, rather than a quantitative human baseline, we conducted a qualitative human expert evaluation of our pipeline's False Positives to understand the boundaries of our cross-encoder's reasoning. 

A manual review of the 3-Gate Max-Gap rejections revealed two primary failure modes:
1. \textbf{Implicit Jurisprudential Shifts:} In $\sim$40\% of errors, the entailing paragraph relied on unstated, assumed knowledge of common law that was not explicitly present in the text, requiring external human legal training to connect the concepts.
2. \textbf{Plausible Alternatives:} In $\sim$60\% of false positives, our model selected a paragraph that provided a logically sound entailment to the query, but simply wasn't the \textit{specific} paragraph the original judge chose to cite. Human reviewers noted that these ``errors'' were often legally defensible, highlighting the extreme difficulty of the task and aligning with established findings that false positives in legal entailment frequently represent plausible, alternative jurisprudential logic \cite{zheng2021when, zhang2025do}. 

Our qualitative analysis proves the inherent challenge of strict historical mapping, further supporting the superiority of tightly bounded, discriminative architectures in preserving precision.

\rev{\subsection{Limitations}

We acknowledge several limitations of our experimental comparison. More recent architectures such as ColBERTv2 \cite{santhanam2022colbertv2}, Longformer \cite{beltagy2020longformer}, and domain-adapted long-context encoders represent important alternative approaches to legal entailment. Our constrained compute environment (a single NVIDIA Tesla T4 GPU in a cloud setting) precluded evaluation of these models within the competition timeline. Additionally, the generative baseline (4-bit quantized Qwen2.5-7B-Instruct) represents a lower bound on current LLM capabilities; larger or more capable generative models may narrow the precision gap we report. Future work should investigate whether long-context encoders can match the precision gains of our decoupled approach without the computational overhead of candidate-level cross-encoding.

% [A6] Missing baselines acknowledgment added per Reviewer 1 feedback

We note that the official COLIEE 2026 Task 2 leaderboard includes several concurrent systems achieving competitive or superior F1 scores. A detailed architectural comparison with these systems is beyond the scope of this paper, as their system descriptions were not available at the time of writing. We defer such cross-system analysis to future post-workshop publications.}

% [A9] Concurrent systems acknowledgment added per Reviewer 2 feedback

\section{Conclusion}

Legal case entailment remains one of the most notoriously difficult tasks in domain-specific natural language processing. The dense, boilerplate-heavy nature of judicial opinions naturally confounds standard retrieval systems, creating Lexical Traps that obscure true semantic causality. In this paper, we introduced a highly specialized, 3-stage Decoupled Architecture designed specifically to untangle this complexity for COLIEE Task 2. By explicitly discarding surrounding context to prevent truncation damage (TDI), and routing the retrieved candidates through a Focal Loss-optimized PMA cross-encoder and a deterministic 3-Gate variance reduction algorithm, we established a high-precision logical boundary that actively resists shallow vocabulary overlap. 

Our empirical ablation suggests that parameter scaling alone is not sufficient for strict logical entailment. On the COLIEE 2026 Task 2 test set, our fine-tuned LegalBERT cross-encoder outperforms a 7B generative model by +0.10 precision, indicating that discriminative boundary detection may be better suited to high-precision legal entailment than autoregressive reasoning under comparable computational constraints. Additionally, our decoupled cross-encoder operates at a fraction of the computational latency required for autoregressive reasoning chains.

Looking forward, we believe the diagnostic frameworks introduced in this study may have applications beyond legal retrieval. \rev{We hypothesize that the LTD framework may have potential applications beyond legal retrieval as a diagnostic signal for evaluating the susceptibility of generative models to lexical collapse across other high-stakes, multi-document domains}, such as medical diagnostics, financial auditing, and regulatory compliance. By providing a framework for detecting when a model's confidence correlates with shallow keyword overlap rather than deep semantic reasoning, this diagnostic may contribute to building more robust AI systems in domains where precision is paramount.

\rev{While our gating algorithm operates at a more conservative recall point, we argue this reflects the correct inductive bias for legal applications where false-positive citations carry significant downstream risk.}

% [A1] Added conclusion recall acknowledgment
% [A2] Softened LTD future work claims

\bibliographystyle{ACM-Reference-Format} 
\bibliography{references}

\end{document}